\documentclass[11pt]{article}
\usepackage[margin=25mm]{geometry}
\usepackage{fontspec}
\setmainfont{Latin Modern Roman}
\usepackage{amsmath,amssymb,mathtools}
\usepackage{unicode-math}
\setmathfont{Latin Modern Math}
\usepackage{xurl,hyperref}
\hypersetup{colorlinks=true,urlcolor=blue}
\setcounter{secnumdepth}{0}
\setlength{\parindent}{0pt}
\setlength{\parskip}{5pt}
\setlength{\emergencystretch}{3em}
\newcommand{\sref}[2]{\hyperlink{source:#1:#2}{[S#2]}}
\newcommand{\cataloguescope}[1]{\paragraph{Recorded target.}#1}
\newcommand{\R}{\mathbb R}
\newcommand{\C}{\mathbb C}
\newcommand{\N}{\mathbb N}
\newcommand{\Z}{\mathbb Z}
\newcommand{\Q}{\mathbb Q}
\newcommand{\F}{\mathbb F}
\newcommand{\eps}{\varepsilon}
\begin{document}
\section{Topology and character sectors in monopole L2-cohomology}\label{statement}

\subsection{Definitions and mathematical statement}

\paragraph{Shared notation.}$\mathbb N=\{1,2,\ldots\}$; the standard number fields have their usual meanings. All additional conventions are stated locally.


\paragraph{Definitions and notation.}For a Riemannian manifold $(M,g)$, $L^2$ norms use the metric volume and pointwise metric on forms. Reduced $L^2$ cohomology is $\ker d/\overline{\operatorname{im}d}$ for the specified closed $L^2$ complex. On a complete manifold its harmonic representatives satisfy $d\alpha=d^*\alpha=0$ and $\int_M|\alpha|^2<\infty$. Higher-dimensional monopole moduli are complete noncompact spaces obtained from gauge equations with specified charges and behavior at infinity. Special holonomy constrains parallel tensors and curvature but does not itself characterize the ends. A topological description of $L^2$ cohomology must account for the metric growth and end conditions. L-modules in Saper’s framework are a proposed tool, not assumed to apply automatically. Extended $L^2$ harmonic forms in the final related question need not themselves be $L^2$; the approximation and residual condition are part of their definition.

\paragraph{Catalogue wording retained for traceability.}
Find a topological interpretation of $L^2$ cohomology for higher-dimensional monopole moduli spaces. Identify suitable tools and decide whether Saper’s L-module techniques can be used. For the associated complete special-holonomy manifolds, understand the asymptotic geometry and $L^2$ cohomology, and find general vanishing theorems, potentially before a complete asymptotic description is available. In the $SU(2)$ monopole case over $\mathbb R^3$, prove that the harmonic forms obtained by the source's cited Segal--Selby construction exhaust all $L^2$ harmonic forms, rather than merely constructing some of them. The following related question in the source asks to study extended $L^2$ harmonic forms on noncompact spaces: forms $\omega$ with $(d+d^*)\omega=0$, admitting locally convergent approximants $\omega_j\in C_0^\infty$ such that $(d+d^*)\omega_j\to0$ in $L^2$.
\paragraph{Scope and source qualification.}The title names character sectors, but the full source asks about the entire moduli-space cohomology and geometric programme. Source remarks about known SU(2) constructions are historical context and do not establish surjectivity here. The surjectivity question in the source remarks and the following anonymous extended-form question are recorded explicitly; the latter is not attributed to Hunsicker. The catalogue title/status was paired with a machine-generated research attempt; that attempt is not peer-reviewed evidence of settlement. The statement covers the original source questions rather than only the title’s subcase.
\subsection{Short English statement}
Give topological descriptions, asymptotics and vanishing results for L² cohomology of monopole and special-holonomy moduli spaces.
\subsection{Research attempt}
\paragraph{Model, date and investigated scope.}
OpenAI GPT-6 Astra (configured model \texttt{gpt-6-astra}), ultra reasoning; generated 11 October 2026. Constructed a non-L2 extended harmonic form with explicit approximants and demonstrated degree-zero metric sensitivity. Monopole exhaustion, L-module/topological interpretation and vanishing results remain unproved.

\paragraph{Sources and literature check.}
Read the supplied complete definition, catalogue wording, scope qualifications and source list for this canonical ID. The original-source citations are inherited as supplied references; no fresh primary-source or novelty search was performed for this attempt. No current open-status claim is inferred from this calculation.

\paragraph{An exact extended-form example.}On the complete line \(\mathbb R\), the constant function \(\omega=1\) satisfies \((d+d^*)\omega=0\) but is not in \(L^2\). Choose a smooth cutoff \(\chi\) with compact support and \(\chi=1\) on \([-1,1]\), and put \(\omega_R(x)=\chi(x/R)\). Then \(\omega_R\to1\) locally smoothly and
\[\|(d+d^*)\omega_R\|_2^2=R^{-1}\int_{\mathbb R}|\chi'(s)|^2ds\longrightarrow0.
\]
Thus the extended-form definition in the source genuinely includes forms outside ordinary L2 cohomology. In \(\mathbb R^d\), the same radial-scale cutoff has gradient norm squared proportional to \(R^{d-2}\), so this particular construction works for \(d=1\) and does not automatically work in higher dimensions. No nonexistence conclusion is drawn from failure of one cutoff choice.
\paragraph{Metric sensitivity of degree zero.}For any connected complete manifold, a square-integrable harmonic zero-form that is both closed and coclosed is constant, since \(df=0\). It belongs to L2 exactly when the total volume is finite (apart from the zero function). Thus even the degree-zero harmonic space can distinguish two complete metrics on the same underlying topology through volume behavior. Topology alone cannot describe L2 cohomology without end hypotheses.
\paragraph{First gap and full scope.}The monopole moduli spaces require their own end estimates and domains. No exhaustion theorem for the source's Segal--Selby forms, L-module applicability, higher-dimensional vanishing theorem or full topological interpretation is proved. Showing that every harmonic form lies in the known constructed sectors needs an upper bound or completeness argument, not merely more examples. This is the first missing step; the extended-form question is retained as the related source question, with no reassignment of its attribution.

\paragraph{Verification and adversarial audit.}
Checked cutoff residual scaling, local convergence, dimension dependence and finite-volume condition for constant L2 forms. This was an LLM adversarial audit, not independent human verification; no theorem prover was used. Prompts A then B applied to this ID: Checked cutoff residual scaling, local convergence, dimension dependence and finite-volume condition for constant L2 forms. Scope remains limited: Prove end/domain estimates giving upper bounds and exhaustion of harmonic forms on the actual monopole moduli spaces.

Coordinator checked the cutoff Dirac norm and dimension scaling, as well as the constant degree-zero volume condition. This additional check was model-assisted, not independent human verification.

\paragraph{Tools, process and contributions.}
Python 3 and local shell used to read supplied files and serialize this manuscript; mathematics derived symbolically in the text, with no numerical experiment or formal proof assistant. Contribution prompt A was followed by prompt B. GPT-6 Astra with ultra reasoning drafted the research and performed its stated self-audit; any additional coordinator contribution is identified in the verification and audit. A GPT-6-based Codex coordinator assembled the manuscript and checked repository integration; its exact variant and reasoning configuration were not exposed. No human mathematical contribution, independent human proof check or formal proof certification is supplied. Complete derivations and the consulted source references are included here; private reasoning transcripts are not published.

\paragraph{Final status and first remaining gap.}
\textbf{UNRESOLVED}. Prove end/domain estimates giving upper bounds and exhaustion of harmonic forms on the actual monopole moduli spaces.

\section{Completion Estimate}
\noindent\textbf{COMPLETION: 0\%}
The calculation is an elementary benchmark or conditional reduction; it establishes no claimed new portion of the complete source target.

\subsection{Sources}
\begin{itemize}
\item[\hypertarget{source:statement:1}{S1}] ULAM.AI and the UnsolvedMath Contributors, UnsolvedMath, catalogue entry. Catalogue curation under CC BY 4.0; definitions and English exposition added locally.
\url{https://huggingface.co/datasets/ulamai/UnsolvedMath}.
\item[\hypertarget{source:statement:2}{S2}] American Institute of Mathematics, $L^2$ harmonic forms in geometry and string theory; notes by Anda Degeratu and Mark Haskins, March 2004, 21 March 2004, Hunsicker questions.
\url{https://aimath.org/WWN/l2harmonic/l2harmonic.pdf}.
\end{itemize}
\label{end:statement}
\end{document}
